\chapter{Evaluation}
\label{chap:evaluation}

The preceding chapters all use event-level F1 (\Cref{chap:classicalml}), but not always under
exactly the same threshold-selection rule. This short chapter names the differences explicitly,
so the comparisons in this report are read correctly rather than taken at face value.

\section{Smoothing and hysteresis}

A second, optional detection path exists, applied with no retraining: the per-window
probabilities are first smoothed with a centered rolling median, then thresholded with
\emph{hysteresis} --- a higher threshold is needed to start a detected event, and a lower
threshold is enough to let an already-started event continue. This generally produces fewer,
cleaner intervals than thresholding the raw probabilities directly. It is applied automatically,
and reported alongside the plain result, for the CNN evaluations in \Cref{chap:deeplearning}; it
is \emph{not} applied to the 12 classical models' headline numbers in
\Cref{tab:classical-results}, which use the plain, unsmoothed path only. This is a real
methodological difference between the two chapters' numbers, named here so it is not mistaken
for an oversight.

\section{Three ways of choosing a threshold}

The ``best threshold'' reported for every classical model, including the augmentation sweep in
\Cref{chap:augmentation}, is chosen and scored on the same out-of-fold predictions, which makes
every classical headline number --- the original 0.894 and the new best, 0.901 --- slightly
optimistic. The ensemble results in \Cref{chap:ensemble} and the CNN results in
\Cref{chap:deeplearning} are more careful: each outer fold's threshold is chosen only from data
outside that fold (other folds, for the ensembles; held-out training recordings, for the CNNs),
never from the fold being scored.

A fully nested re-scoring of the classical models, under this same stricter rule, is not
finished --- only 3 models have been run this way so far (the full run is memory-heavy and was
interrupted once for lack of RAM), and none of them is the new best model (SVM, channel
dropout). Until it is, the comparison between the classical models' 0.894/0.901 and the best
CNN's 0.817 should be read as informative but not perfectly matched: the classical numbers carry
a small optimism the CNN numbers do not.

\section{Where this leaves the comparison}

\begin{table}[H]
  \centering
  \caption{The project's results, side by side, with the threshold protocol each one used.}
  \label{tab:final-comparison}
  \small
  \begin{tabular}{p{7.3cm}rl}
    \toprule
    Result & Event F1 & Threshold protocol \\
    \midrule
    \textbf{SVM, channel-dropout augmentation (best overall)} & \textbf{0.901} & same-fold (optimistic) \\
    SVM + ShallowConvNet ensemble (ties the best) & 0.901 & held-out-fold \\
    SVM, unbalanced (original baseline) & 0.894 & same-fold (optimistic) \\
    SVM + Random Forest ensemble (ties the original baseline) & 0.894 & held-out-fold \\
    ShallowConvNet, amplitude alone (best CNN) & 0.817 & nested \\
    \bottomrule
  \end{tabular}
\end{table}

Taken together: a single classical model, SVM trained with channel-dropout augmentation, is both
the simplest and the best-performing result in this project. Every attempt to do better by
blending models ties it at best, never beats it. The CNN, trained end-to-end on raw signal, does
not come close to the feature-engineered approach at this dataset size, no matter which of its 8
augmentation transforms is used --- exactly the outcome the project set out to test for.

\begin{summarybox}
Three threshold-selection rules are in play across this report --- same-fold (slightly
optimistic), held-out-fold, and nested --- and they are not interchangeable. The one piece of
evaluation work still missing is a full nested re-scoring of the classical models (including the
new best one), which would make the classical-vs-CNN comparison fully fair; until then, the
classical numbers should be read as slightly favorable to the classical side relative to the
CNNs' nested 0.817. Even allowing for that gap, the margin (0.901 vs.\ 0.817) is wide enough that
the conclusion --- classical beats deep learning here --- is not in doubt.
\end{summarybox}
